\chapter{External Distribution Sort, Sovracampionamento e Campionamento Casuale}
\label{chap:external-distribution-sort}

\begin{center}
  \textbf{Lezione 6}\\
  \emph{External Distribution Sort, Over-Sampling con Limiti di Chernoff e Campionamento Casuale}\\[0.5em]
  Data: 30 Settembre 2026
\end{center}

Nel contesto dei modelli di memoria gerarchica e dell'ordinamento per memorie di massa (\textit{External Memory Model}), il \textbf{Multiway Merge Sort} studiato nel \autoref{cha:lecture-three} rappresenta l'approccio \textit{bottom-up} canonico: i dati vengono inizialmente suddivisi in sequenze ordinate (\textit{runs}) e successivamente fusi ricorsivamente a $k$ vie fino a ottenere un unico vettore ordinato.

Tuttavia, esiste un paradigma duale di fondamentale rilevanza teorica e applicativa: l'\textbf{External Distribution Sort} (talvolta denominato \textit{Multiway Quicksort} esterno). Tale approccio adotta una filosofia \textit{top-down}: invece di fondere blocchi già ordinati, distribuisce l'intero insieme di dati non ordinati in molteplici partizioni disgiunte (\textit{buckets}), ciascuna associata a un intervallo ordinato di valori di chiave, per poi procedere ricorsivamente su ciascun secchio in modo indipendente.

In questo capitolo esaminiamo l'architettura dettagliata dell'External Distribution Sort, affrontando la sfida critica della selezione bilanciata dei pivot mediante la tecnica statistica dell'\textbf{Over-Sampling} (sovracampionamento). Analizziamo quindi formalmente il limite di fallimento mediante i \textbf{Chernoff Bounds} (disuguaglianze di concentrazione della misura) e approfondiamo gli algoritmi di \textbf{Campionamento Casuale Uniforme} (\textit{Random Sampling}), confrontando i metodi basati su tabelle hash con il paradigma streaming del \textbf{Reservoir Sampling}.

\section{External Distribution Sort e Partizionamento a \texorpdfstring{$k$}{k} Vie}
\label{sec:external-distribution-sort-model}

Consideriamo un dataset costituito da $n$ elementi memorizzati su disco secondario nel modello a due livelli di Aggarwal e Vitter~\cite{aggarwal1988}, caratterizzato da una memoria centrale interna di capacità pari a $M$ elementi e da trasferimenti a blocchi di dimensione $B$ elementi. L'obiettivo consiste nell'ordinare l'array $S[1 \dots n]$ minimizzando il numero complessivo di operazioni di I/O.

\subsection{Meccanismo Operativo e Definizione dei Bucket}
L'algoritmo suddivide l'array $S$ in $k$ partizioni continue e disgiunte, denominate \textbf{secchi} (\textit{buckets}) $B_1, B_2, \dots, B_k$, mediante la selezione di $k - 1$ perni (\textit{pivots}) ordinati in modo strettamente crescente:
\begin{equation}
  s_1 < s_2 < \dots < s_{k-1}.
\end{equation}

Per garantire una trattazione uniforme dei casi di frontiera, si introducono due valori sentinella concettuali:
\begin{equation}
  s_0 \coloneqq -\infty, \qquad s_k \coloneqq +\infty.
\end{equation}

\begin{definition}[Secchio o Bucket $B_i$]
Per ogni indice $i \in \{1, \dots, k\}$, il secchio $B_i$ è il sottoinsieme degli elementi di $S$ le cui chiavi cadono nell'intervallo semi-aperto delimitato dai pivot adiacenti:
\begin{equation}
  B_i \coloneqq \left\{ x \in S \mid s_{i-1} < x \le s_i \right\}.
  \label{eq:bucket-definition}
\end{equation}
\end{definition}

La topologia del partizionamento a $k$ vie è schematizzata nella \autoref{fig:k-way-partitioning-model}.

\begin{figure}[H]
  \centering
  \resizebox{0.95\textwidth}{!}{%
  \begin{tikzpicture}[
    scale=1.0,
    bucket/.style={rectangle, draw=blue!80!black, thick, minimum height=1.3cm, align=center, font=\small\bfseries},
    sentinel/.style={font=\footnotesize\bfseries, text=red!80!black},
    pivot/.style={font=\footnotesize\bfseries, text=blue!85!black},
    ptr/.style={-{Stealth[scale=0.9]}, very thick, draw=blue!85!black}
  ]
    % Segmenti dei bucket
    \node[bucket, fill=blue!12, minimum width=2.8cm, anchor=west] (b1) at (0, 0) {
      $B_1$\\[2pt]\footnotesize ($s_0 < x \le s_1$)
    };
    \node[bucket, fill=orange!15, minimum width=2.8cm, anchor=west] (b2) at (b1.east) {
      $B_2$\\[2pt]\footnotesize ($s_1 < x \le s_2$)
    };
    \node[bucket, fill=teal!15, minimum width=2.8cm, anchor=west] (b3) at (b2.east) {
      $B_3$\\[2pt]\footnotesize ($s_2 < x \le s_3$)
    };
    \node[draw=none, minimum width=1.5cm, anchor=west, align=center, font=\Large\bfseries] (dots) at (b3.east) {
      $\dots$
    };
    \node[bucket, fill=purple!15, minimum width=2.8cm, anchor=west] (bk) at (dots.east) {
      $B_k$\\[2pt]\footnotesize ($s_{k-1} < x \le s_k$)
    };

    % Pivot e Sentinelle in alto
    \node[above=14pt, sentinel] at (b1.north west) {$s_0 = -\infty$};
    \draw[ptr] ([yshift=18pt]b1.north east) -- (b1.north east)
      node[pos=0.0, above, pivot] {$s_1$};
    \draw[ptr] ([yshift=18pt]b2.north east) -- (b2.north east)
      node[pos=0.0, above, pivot] {$s_2$};
    \draw[ptr] ([yshift=18pt]b3.north east) -- (b3.north east)
      node[pos=0.0, above, pivot] {$s_3$};
    \draw[ptr] ([yshift=18pt]bk.north west) -- (bk.north west)
      node[pos=0.0, above, pivot] {$s_{k-1}$};
    \node[above=14pt, sentinel] at (bk.north east) {$s_k = +\infty$};

    % Graffe decorative inferiori
    \draw[decorate, decoration={brace, amplitude=5pt, mirror}, thick, blue!85!black]
      ([yshift=-4pt]b1.south west) -- ([yshift=-4pt]b1.south east)
      node[midway, below=7pt, font=\scriptsize\bfseries] {$|B_1| \approx \frac{n}{k}$};

    \draw[decorate, decoration={brace, amplitude=5pt, mirror}, thick, orange!90!black]
      ([yshift=-4pt]b2.south west) -- ([yshift=-4pt]b2.south east)
      node[midway, below=7pt, font=\scriptsize\bfseries] {$|B_2| \approx \frac{n}{k}$};

    \draw[decorate, decoration={brace, amplitude=5pt, mirror}, thick, teal!90!black]
      ([yshift=-4pt]b3.south west) -- ([yshift=-4pt]b3.south east)
      node[midway, below=7pt, font=\scriptsize\bfseries] {$|B_3| \approx \frac{n}{k}$};

    \draw[decorate, decoration={brace, amplitude=5pt, mirror}, thick, purple!85!black]
      ([yshift=-4pt]bk.south west) -- ([yshift=-4pt]bk.south east)
      node[midway, below=7pt, font=\scriptsize\bfseries] {$|B_k| \approx \frac{n}{k}$};
  \end{tikzpicture}%
  }
  \caption{Partizionamento dell'universo delle chiavi in $k$ secchi disgiunti mediante $k-1$ pivot ordinati e due sentinelle concettuali.}
  \label{fig:k-way-partitioning-model}
\end{figure}

\begin{property}[Proprietà di Concatenazione Immediata]
A differenza del Merge Sort, in cui la fase conclusiva richiede la complessa fusione di flussi ordinati, l'External Distribution Sort gode della proprietà per cui:
\begin{equation}
  \forall x \in B_i, \quad \forall y \in B_j \quad \text{con } i < j \implies x \le y.
\end{equation}
Pertanto, una volta che ciascun singolo secchio $B_i$ è stato ricorsivamente ordinato, l'array globale ordinato coincide semplicemente con la concatenazione sequenziale diretta dei secchi:
\begin{equation}
  S_{\text{ordinato}} = B_1 \circ B_2 \circ \dots \circ B_k.
\end{equation}
Non è necessaria alcuna operazione di fusione a posteriori: il risparmio di computazione e di gestione puntatori nella fase di ricomposizione è integrale.
\end{property}

\subsection{L'Obiettivo di Bilanciamento e l'Impatto sull'I/O}
Affinché l'algoritmo consegua la complessità ottima di I/O nel modello esterno, la scelta del numero di rami $k$ e delle posizioni dei pivot deve rispettare due requisiti architetturali stringenti:

\begin{enumerate}
  \item \textbf{Capacità dei buffer interni:} Durante la scansione di $S$, per ciascuno dei $k$ bucket deve essere mantenuto in memoria rapida almeno un blocco di scrittura di dimensione $B$ (\textit{Write Buffer}). Poiché la memoria interna ha capacità $M$, il grado massimo di partizionamento ammissibile è:
  \begin{equation}
    k = \Theta\left(\frac{M}{B}\right), \quad \text{con } k \ge 2.
    \label{eq:k-branching-factor}
  \end{equation}
  Quando il buffer associato al bucket $B_i$ si riempie, esso viene scaricato su disco con una singola scrittura I/O di $B$ elementi, realizzando un flusso strettamente sequenziale.
  
  \item \textbf{Bilanciamento dimensionale dei secchi:} La dimensione di ciascun bucket deve risultare strettamente uniforme:
  \begin{equation}
    |B_i| = \Theta\left(\frac{n}{k}\right), \quad \forall i \in \{1, \dots, k\}.
    \label{eq:balance-condition}
  \end{equation}
\end{enumerate}

\begin{theorem}[Complessità I/O dell'External Distribution Sort Bilanciato]
Se la condizione di bilanciamento~\eqref{eq:balance-condition} è soddisfatta ad ogni passo della ricorsione, l'altezza complessiva dell'albero di distribuzione è pari a:
\begin{equation}
  h = \mathcal{O}\left(\log_k \frac{n}{M}\right) = \mathcal{O}\left(\log_{M/B} \frac{n}{B}\right).
\end{equation}
Poiché ciascun livello di partizionamento richiede la scansione lineare dell'intero dataset con costo $\mathcal{O}(n/B)$ trasferimenti di blocco, la complessità totale di I/O dell'algoritmo è:
\begin{equation}
  \text{I/O}(n) = \mathcal{O}\left(\frac{n}{B} \log_{M/B} \frac{n}{B}\right),
\end{equation}
raggiungendo esattamente il limite inferiore asintotico di ordinamento stabilito dal Teorema di Aggarwal e Vitter.
\end{theorem}

\begin{property}[Rischio di Degenerazione da Sbilanciamento]
Qualora anche un singolo secchio $B_j$ risulti sovradimensionato (ad esempio $|B_j| = \Omega(n)$), l'altezza dell'albero ricorsivo non decresce a passo logaritmico, ma degrada linearmente verso $\Omega(n/M)$. In tale scenario avverso, il costo di I/O collasserebbe verso una complessità quadratica $\mathcal{O}((n/B)^2)$, rendendo l'algoritmo inutilizzabile su dataset massivi. La corretta scelta dei pivot costituisce pertanto l'asse portante dell'intero algoritmo.
\end{property}

\section{La Tecnica dell'Over-Sampling (Sovracampionamento)}
\label{sec:oversampling-technique}

Come determinare $k - 1$ pivot in grado di garantire con elevatissima confidenza che ciascuno dei $k$ secchi contenga $\Theta(n/k)$ elementi?

Un approccio ingenuo consisterebbe nell'estrarre casualmente esattamente $k - 1$ elementi da $S$ e ordinarli. Tuttavia, dalla teoria delle statistiche d'ordine per variabili casuali uniformi è noto che la dimensione del secchio massimo generato da $k-1$ campioni indipendenti presenta una varianza elevata, con una taglia attesa del secchio massimo pari a:
\begin{equation}
  \mathbb{E}\left[\max_{1 \le i \le k} |B_i|\right] \approx \frac{n}{k} \ln k.
\end{equation}
Per valori realistici di $k$ (ad esempio $k = M/B \approx 10^3 \div 10^5$), il fattore moltiplicativo $\ln k \approx 7 \div 12$ comporta uno sbilanciamento intollerabile.

Per risolvere radicalmente questa criticità senza dover calcolare i quantili esatti (operazione che richiederebbe costosi algoritmi di selezione lineare su disco), si impiega la tecnica del \textbf{Sovracampionamento} (\textit{Over-Sampling}).

\subsection{I Parametri e la Procedura Operativa}
Si introduce un moltiplicatore di campionamento $a = \Theta(\log_2 k)$, denominato \textbf{fattore di over-sampling}. Nei dettagli analitici dimostrati nel seguito, poniamo formalmente:
\begin{equation}
  a \coloneqq 12 \ln k.
  \label{eq:oversampling-factor}
\end{equation}

La procedura operativa di selezione si articola in tre fasi elementari:

\begin{enumerate}
  \item \textbf{Estrazione del campione globale:} Si estraggono in modo uniformemente casuale e indipendente dall'array $S$ un numero totale di campioni pari a:
  \begin{equation}
    |A| = (a + 1)k - 1.
    \label{eq:sample-size}
  \end{equation}
  Tali elementi vengono memorizzati in un buffer ausiliario $A$.
  
  \item \textbf{Ordinamento interno del campione:} L'array $A$ viene ordinato interamente in memoria centrale mediante un algoritmo ottimale in tempo CPU (ad esempio Quicksort o Merge Sort). Poiché la cardinalità $|A| = \mathcal{O}(k \ln k) = \mathcal{O}(\frac{M}{B} \ln \frac{M}{B})$ è di ordini di grandezza inferiore alla capienza totale della RAM $M$, l'ordinamento di $A$ avviene con zero trasferimenti su disco (0 I/O) e richiede una frazione trascurabile di millisecondo.
  
  \item \textbf{Selezione equidistante dei Pivot:} I $k - 1$ pivot vengono selezionati estraendo dall'array ordinato $A$ gli elementi posti a passo regolare $(a + 1)$:
  \begin{equation}
    s_i \coloneqq A[(a + 1) \cdot i], \quad \text{per } i = 1, \dots, k - 1.
    \label{eq:pivot-selection-rule}
  \end{equation}
\end{enumerate}

La struttura interna del campione ordinato $A$ è illustrata dettagliatamente nella \autoref{fig:oversampling-sample-layout}.

\begin{figure}[H]
  \centering
  \resizebox{0.95\textwidth}{!}{%
  \begin{tikzpicture}[
    scale=1.0,
    box/.style={rectangle, draw=black!80, thick, minimum height=1.1cm, align=center, font=\small},
    pivotBox/.style={rectangle, draw=blue!90!black, fill=blue!20, thick, minimum height=1.1cm, minimum width=1.1cm, align=center, font=\small\bfseries},
    arr/.style={-{Stealth[scale=0.9]}, very thick, draw=blue!85!black}
  ]
    % Celle del campione A
    \node[box, fill=gray!12, minimum width=2.4cm, anchor=west] (gap0) at (0, 0) {$a$ campioni};
    \node[pivotBox, anchor=west] (p1) at (gap0.east) {$s_1$};
    
    \node[box, fill=gray!12, minimum width=2.4cm, anchor=west] (gap1) at (p1.east) {$a$ campioni};
    \node[pivotBox, anchor=west] (p2) at (gap1.east) {$s_2$};
    
    \node[box, fill=gray!12, minimum width=2.4cm, anchor=west] (gap2) at (p2.east) {$a$ campioni};
    \node[draw=none, minimum width=1.2cm, anchor=west, font=\large\bfseries] (dots) at (gap2.east) {$\dots$};
    \node[pivotBox, anchor=west] (pk) at (dots.east) {$s_{k-1}$};
    
    \node[box, fill=gray!12, minimum width=2.4cm, anchor=west] (gapk) at (pk.east) {$a$ campioni};

    % Indici di riga del vettore A
    \draw[arr] ([yshift=16pt]p1.north) -- (p1.north)
      node[pos=0.0, above, font=\scriptsize\bfseries, text=blue!90!black] {$A[a+1]$};
      
    \draw[arr] ([yshift=16pt]p2.north) -- (p2.north)
      node[pos=0.0, above, font=\scriptsize\bfseries, text=blue!90!black] {$A[2(a+1)]$};

    \draw[arr] ([yshift=16pt]pk.north) -- (pk.north)
      node[pos=0.0, above, font=\scriptsize\bfseries, text=blue!90!black] {$A[(k-1)(a+1)]$};

    % Graffa inferiore totale
    \draw[decorate, decoration={brace, amplitude=6pt, mirror}, thick]
      ([yshift=-4pt]gap0.south west) -- ([yshift=-4pt]gapk.south east)
      node[midway, below=8pt, font=\small\bfseries] {Campione globale ordinato $|A| = (a+1)k - 1$};
  \end{tikzpicture}%
  }
  \caption{Topologia dell'array ausiliario ordinato $A$: tra ciascuna coppia di pivot contigui cadono esattamente $a$ campioni intermedi.}
  \label{fig:oversampling-sample-layout}
\end{figure}

\begin{property}[Invariante dei Campioni Intermedi]
Per costruzione dell'algoritmo, tra due pivot consecutivi $s_{i-1}$ e $s_i$ cadono \textbf{esattamente $a$ campioni} estratti e presenti in $A$. Poiché vi sono $k - 1$ pivot, si hanno $k$ intervalli intermedi, ciascuno popolato da $a$ campioni, per un totale di $k \cdot a$ campioni non-pivot. Sommando i $k - 1$ pivot stessi, la cardinalità totale risulta:
\begin{equation}
  k \cdot a + (k - 1) = (a + 1)k - 1 = |A|,
\end{equation}
confermando la perfetta conservazione della taglia del campione.
\end{property}

\subsection{Esempio Numerico Tracciato di Selezione}
\begin{example}[Traccia di Sovracampionamento con $n=64$, $k=4$, $a=3$]
Per comprendere concretamente l'algoritmo, consideriamo un insieme didattico con $n = 64$ elementi, per il quale vogliamo generare $k = 4$ partizioni. Fissiamo il fattore di sovracampionamento a un valore didattico $a = 3$:
\begin{enumerate}
  \item Dimensione del campione estratto:
  \begin{equation}
    |A| = (a + 1)k - 1 = (3 + 1) \cdot 4 - 1 = 15 \text{ elementi}.
  \end{equation}
  \item Supponiamo che l'estrazione casuale e il successivo ordinamento in memoria di $A$ producano il vettore ordinato:
  \begin{equation*}
    A = [\, 4,\, 7,\, 9,\, \mathbf{12},\, 15,\, 18,\, 22,\, \mathbf{29},\, 33,\, 38,\, 41,\, \mathbf{50},\, 55,\, 61,\, 68 \,].
  \end{equation*}
  \item Il passo regolare di selezione è $a + 1 = 4$. I $k - 1 = 3$ pivot sono quindi:
  \begin{align*}
    s_1 &= A[1 \cdot 4] = A[4] = \mathbf{12}, \\
    s_2 &= A[2 \cdot 4] = A[8] = \mathbf{29}, \\
    s_3 &= A[3 \cdot 4] = A[12] = \mathbf{50}.
  \end{align*}
  \item I 4 secchi risultanti per la successiva fase di partizionamento saranno:
  \begin{align*}
    B_1 &= \{ x \in S \mid -\infty < x \le 12 \}, \\
    B_2 &= \{ x \in S \mid 12 < x \le 29 \}, \\
    B_3 &= \{ x \in S \mid 29 < x \le 50 \}, \\
    B_4 &= \{ x \in S \mid 50 < x \le +\infty \}.
  \end{align*}
\end{enumerate}
Ciascun intervallo $(s_{i-1}, s_i]$ racchiude esattamente $a = 3$ elementi del campione $A$, garantendo che la densità stimata delle chiavi rimanga proporzionata alla distribuzione empirica dei dati.
\end{example}

\section{Dimostrazione del Limite di Fallimento via Chernoff Bound}
\label{sec:chernoff-bound-proof}

Vogliamo ora dimostrare formalmente che l'estrazione casuale tramite over-sampling garantisce il bilanciamento dimensionale dei secchi con probabilità schiacciante.

\subsection{Definizione Formale di Fallimento}
Diciamo che l'estrazione dei pivot si risolve in un \textbf{fallimento} se almeno uno dei $k$ secchi $B_i$ risulta gravemente sbilanciato, contenendo un numero di elementi pari o superiore al quadruplo della taglia media attesa $n/k$:
\begin{equation}
  \mathbb{P}(\text{fallimento}) \coloneqq \mathbb{P}\left(\exists i \in \{1, \dots, k\} : |B_i| \ge \frac{4n}{k}\right).
  \label{eq:failure-definition}
\end{equation}

\subsection{Costruzione Geometrica sull'Array Ordinato Virtuale \texorpdfstring{$S_{\text{sort}}$}{S\_sort}}
Per analizzare rigorosamente l'evento di fallimento, introduciamo una costruzione geometrica ausiliaria. Si consideri l'array $S$ concettualmente ordinato, denotato con $S_{\text{sort}}$.

Suddividiamo idealmente l'intero array $S_{\text{sort}}$ in $\frac{k}{2}$ intervalli contigui di ampiezza uniforme, indicati con $t_j$ per $j = 1, \dots, \frac{k}{2}$:
\begin{equation}
  |t_j| \coloneqq \frac{n}{k/2} = \frac{2n}{k}.
  \label{eq:segment-width}
\end{equation}

La relazione geometrica tra i segmenti canonici $t_j$ e un secchio sovradimensionato $B_i$ è visualizzata nella \autoref{fig:pigeonhole-chernoff-construction}.

\begin{figure}[H]
  \centering
  \resizebox{0.95\textwidth}{!}{%
  \begin{tikzpicture}[
    scale=1.0,
    segment/.style={rectangle, draw=black!80, thick, minimum height=1.1cm, align=center, font=\small},
    bucketBox/.style={rectangle, draw=red!80!black, fill=red!15, very thick, minimum height=1.2cm, align=center, font=\small\bfseries}
  ]
    % Array ordinato S_sort diviso in k/2 segmenti di taglia 2n/k
    \node[segment, fill=gray!10, minimum width=2.4cm, anchor=west] (t1) at (0, 0) {$t_1$\\\footnotesize($\frac{2n}{k}$)};
    \node[segment, fill=yellow!20, minimum width=2.4cm, anchor=west] (t2) at (t1.east) {$t_2$\\\footnotesize($\frac{2n}{k}$)};
    \node[segment, fill=orange!25, minimum width=2.4cm, anchor=west] (t3) at (t2.east) {$t^* = t_3$\\\footnotesize($\frac{2n}{k}$)};
    \node[segment, fill=yellow!20, minimum width=2.4cm, anchor=west] (t4) at (t3.east) {$t_4$\\\footnotesize($\frac{2n}{k}$)};
    \node[segment, fill=gray!10, minimum width=2.4cm, anchor=west] (t5) at (t4.east) {$t_5$\\\footnotesize($\frac{2n}{k}$)};
    \node[draw=none, minimum width=1.2cm, anchor=west, font=\large\bfseries] (dots) at (t5.east) {$\dots$};
    \node[segment, fill=gray!10, minimum width=2.4cm, anchor=west] (tk) at (dots.east) {$t_{k/2}$\\\footnotesize($\frac{2n}{k}$)};

    % Secchio Bi sovradimensionato che copre interamente t3
    \node[bucketBox, minimum width=6.6cm, anchor=south] (bi) at ([yshift=14pt]t3.north) {
      Secchio Sbilanciato $B_i$ con $|B_i| \ge \frac{4n}{k}$\\
      \footnotesize(Contiene interamente $t_3 \implies t_3 \subseteq B_i$)
    };

    \draw[-{Stealth[scale=1.0]}, very thick, red!80!black] (bi.south) -- (t3.north);

    % Delimitatori di S_sort
    \node[above=4pt, font=\scriptsize\bfseries] at (t1.north west) {$S_{\text{sort}}[1]$};
    \node[above=4pt, font=\scriptsize\bfseries] at (tk.north east) {$S_{\text{sort}}[n]$};
  \end{tikzpicture}%
  }
  \caption{Costruzione geometrica di contenimento: un secchio $B_i$ di taglia $\ge 4n/k$ copre interamente almeno un segmento logico $t^* \subset S_{\text{sort}}$ di ampiezza $2n/k$.}
  \label{fig:pigeonhole-chernoff-construction}
\end{figure}

\begin{lemma}[Lemma di Contenimento Geometrizzato]
\label{lem:containment}
Se un secchio $B_i$ possiede cardinalità $|B_i| \ge \frac{4n}{k}$, allora esso \textbf{contiene interamente al suo interno almeno uno} dei segmenti canonici $t_j$:
\begin{equation}
  \exists j \in \left\{1, \dots, \frac{k}{2}\right\} \quad \text{tale che} \quad t_j \subseteq B_i.
\end{equation}
\end{lemma}

\begin{proof}
La dimostrazione segue dal principio dei cassetti continuo. I punti di separazione tra i segmenti consecutivi $t_j$ distano esattamente $\frac{2n}{k}$ l'uno dall'altro. Se un intervallo continuo $B_i$ ha lunghezza pari ad almeno $\frac{4n}{k}$, la sua ampiezza è pari al doppio della distanza tra due delimitatori adiacenti. Di conseguenza, $B_i$ deve necessariamente oltrepassare almeno due delimitatori consecutivi di segmento, inglobando interamente al proprio interno l'intervallo $t_j$ compreso tra essi.
\end{proof}

\subsection{La Contraddizione sui Campioni Estratti}
Mettiamo ora a confronto il contenimento geometrico con la costruzione dell'algoritmo di over-sampling:
\begin{itemize}
  \item Per costruzione, il secchio $B_i$ è delimitato da due pivot consecutivi $s_{i-1}$ e $s_i$.
  \item In virtù della regola di campionamento~\eqref{eq:pivot-selection-rule}, nell'array ordinato $A$ vi sono \textbf{esattamente $a$ campioni} strettamente compresi tra $s_{i-1}$ e $s_i$.
  \item Pertanto, il secchio $B_i$ può contenere al suo interno al più $a$ elementi appartenenti al campione estratto $A$.
  \item Poiché per il \autoref{lem:containment} il segmento $t_j$ è un sottoinsieme di $B_i$ ($t_j \subseteq B_i$), il numero di campioni di $A$ che cadono in $t_j$ non può superare il numero di campioni che cadono in $B_i$.
\end{itemize}

Ne discende la deduzione fondamentale:
\begin{equation}
  \left( |B_i| \ge \frac{4n}{k} \right) \implies \left( \exists j \in \left\{1, \dots, \frac{k}{2}\right\} : |t_j \cap A| < a + 1 \right).
\end{equation}

Possiamo dunque maggiorare la probabilità di fallimento dell'intero algoritmo mediante la probabilità dell'unione di eventi sui segmenti $t_j$:
\begin{align}
  \mathbb{P}(\text{fallimento}) &\le \mathbb{P}\left(\exists j \in \left\{1, \dots, \frac{k}{2}\right\} : t_j \text{ contiene } < a + 1 \text{ campioni di } A\right) \notag \\
  &\le \sum_{j=1}^{k/2} \mathbb{P}\left(t_j \text{ contiene } < a + 1 \text{ campioni di } A\right). \label{eq:union-bound-step}
\end{align}
Dato che tutti i segmenti $t_j$ hanno la medesima lunghezza $\frac{2n}{k}$, la probabilità è identica per ogni $j$. Fissato un generico intervallo target $t^*$, applicando lo \textit{Union Bound}:
\begin{equation}
  \mathbb{P}(\text{fallimento}) \le \frac{k}{2} \cdot \mathbb{P}\left(t^* \text{ contiene } < a + 1 \text{ campioni di } A\right).
  \label{eq:union-bound-final}
\end{equation}

\subsection{Formalizzazione con Variabili Indicatrici e Valore Atteso}
Sia $t^*$ un segmento prefissato di lunghezza $|t^*| = \frac{2n}{k}$. Per ciascun elemento campionato $A[r] \in A$ (per $r = 1, \dots, |A|$), definiamo la variabile casuale bernoulliana:
\begin{equation}
  X_r \coloneqq \begin{cases}
    1 & \text{se il campione } A[r] \in t^*, \\
    0 & \text{altrimenti}.
  \end{cases}
\end{equation}

Poiché i campioni sono estratti uniformemente a caso da $S$, la probabilità di successo per ogni singola estrazione coincide con la frazione dimensionale occupata da $t^*$ rispetto all'intero array $S$:
\begin{equation}
  p \coloneqq \mathbb{P}(X_r = 1) = \frac{|t^*|}{n} = \frac{\frac{2n}{k}}{n} = \frac{2}{k}.
\end{equation}

La variabile casuale che conta il numero complessivo di campioni caduti in $t^*$ è data dalla somma:
\begin{equation}
  X \coloneqq \sum_{r=1}^{|A|} X_r.
\end{equation}

Calcoliamo il valore atteso $\mathbb{E}[X]$ sfruttando la linearità dell'aspettazione:
\begin{equation}
  \mathbb{E}[X] = \sum_{r=1}^{|A|} \mathbb{E}[X_r] = |A| \cdot p = [(a + 1)k - 1] \cdot \frac{2}{k} = 2(a + 1) - \frac{2}{k}.
\end{equation}

Essendo il numero di partizioni $k \ge 2$, il termine sottrattivo soddisfa $\frac{2}{k} \le 1$. Pertanto:
\begin{equation}
  \mathbb{E}[X] \ge 2(a + 1) - 1.
\end{equation}
Poiché $a \ge 1$, si ha banalmente che $1 \le \frac{1}{2}(a + 1)$, da cui discende la stima conservativa:
\begin{equation}
  \mathbb{E}[X] \ge 2(a + 1) - \frac{1}{2}(a + 1) = \frac{3}{2}(a + 1).
  \label{eq:expected-value-bound}
\end{equation}

Invertendo la relazione tra la soglia $a + 1$ e il valore atteso $\mathbb{E}[X]$:
\begin{equation}
  a + 1 \le \frac{2}{3} \mathbb{E}[X] = \left(1 - \frac{1}{3}\right) \mathbb{E}[X].
  \label{eq:delta-derivation}
\end{equation}

\subsection{Applicazione del Chernoff Bound sulla Coda Sinistra}
La formulazione classica del \textbf{Limite Inferiore di Chernoff}~\cite{chernoff1952,motwani1995} stabilisce che per una somma $X = \sum X_r$ di variabili bernoulliane indipendenti e per qualsiasi fattore di deviazione $\delta \in (0, 1)$:
\begin{equation}
  \mathbb{P}(X < (1 - \delta)\mathbb{E}[X]) \le e^{-\frac{\delta^2}{2} \mathbb{E}[X]}.
  \label{eq:chernoff-lower-tail}
\end{equation}

Identificando $\delta = \frac{1}{3}$ in conformità all'\autoref{eq:delta-derivation}:
\begin{equation}
  \mathbb{P}(X < a + 1) \le \mathbb{P}\left(X < \left(1 - \frac{1}{3}\right)\mathbb{E}[X]\right) \le e^{-\frac{(1/3)^2}{2} \mathbb{E}[X]} = e^{-\frac{1}{18} \mathbb{E}[X]}.
\end{equation}

Sostituendo nell'esponente la limitazione inferiore del valore atteso ottenuta nell'\autoref{eq:expected-value-bound} ($\mathbb{E}[X] \ge \frac{3}{2}(a + 1)$):
\begin{equation}
  -\frac{1}{18} \mathbb{E}[X] \le -\frac{1}{18} \cdot \frac{3}{2}(a + 1) = -\frac{a + 1}{12}.
\end{equation}

Introducendo il valore analitico stabilito per il fattore di over-sampling $a = 12 \ln k$:
\begin{equation}
  -\frac{a + 1}{12} < -\frac{a}{12} = -\frac{12 \ln k}{12} = -\ln k.
\end{equation}

Sostituendo all'interno dell'esponenziale:
\begin{equation}
  \mathbb{P}(X < a + 1) \le e^{-\ln k} = \frac{1}{k}.
  \label{eq:single-interval-bound}
\end{equation}

\subsection{Conclusione della Dimostrazione e Strategia Las Vegas}
Possiamo ora combinare la stima del singolo intervallo~\eqref{eq:single-interval-bound} con l'\autoref{eq:union-bound-final} derivata dallo Union Bound su tutti i $\frac{k}{2}$ segmenti candidati:
\begin{equation}
  \mathbb{P}(\text{fallimento}) \le \frac{k}{2} \cdot \mathbb{P}(X < a + 1) \le \frac{k}{2} \cdot \frac{1}{k} = \frac{1}{2}.
  \label{eq:total-failure-bound}
\end{equation}

\begin{theorem}[Garanzia di Successo dell'Over-Sampling]
Con la scelta del parametro $a = 12 \ln k$, la probabilità che \textbf{tutti} i $k$ secchi risultino rigorosamente bilanciati soddisfa:
\begin{equation}
  \mathbb{P}(\text{successo}) = 1 - \mathbb{P}(\text{fallimento}) \ge 1 - \frac{1}{2} = \frac{1}{2} \quad (50\%).
\end{equation}
\end{theorem}

\begin{property}[Ripetizione a Costo Ammortizzato Nullo (\textit{Las Vegas})]
Qualora la verifica dimensionale rilevi che un secchio supera la taglia ammissibile $\frac{4n}{k}$, l'algoritmo rigetta i pivot ed esegue un nuovo campionamento indipendente. Il numero di iterazioni necessarie prima di ottenere un partizionamento ammissibile è una variabile casuale geometrica $T \sim \text{Geom}(p)$ con parametro di successo $p \ge \frac{1}{2}$. 
Il numero atteso di iterazioni è:
\begin{equation}
  \mathbb{E}[T] = \frac{1}{p} \le 2.
\end{equation}
Poiché l'estrazione e l'ordinamento di $A$ richiedono esclusivamente tempo CPU in RAM con 0 I/O su disco, ripetere il campionamento 1 o 2 volte introduce un overhead architetturale trascurabile, preservando con certezza la complessità ottima $\mathcal{O}\left(\frac{n}{B} \log_{M/B}\frac{n}{B}\right)$ I/O.
\end{property}

\section{Algoritmi di Campionamento Casuale Uniforme (\textit{Random Sampling})}
\label{sec:random-sampling-algorithms}

La validità probabilistica dell'External Distribution Sort presuppone che la fase 1 sia in grado di estrarre campioni dall'array $S$ in modo strettamente uniforme e senza reinserimento.

\subsection{Formulazione del Problema}
Dato un insieme discreto di $n$ elementi indicizzati $S[1 \dots n]$, si richiede di estrarre un sottoinsieme $D \subset S$ di cardinalità fissa $|D| = m$ (nel nostro caso $m = |A| = (a+1)k - 1$) in modo tale che ogni possibile combinazione di $m$ elementi distinti abbia la medesima probabilità di estrazione, ossia che:
\begin{equation}
  \mathbb{P}(S[i] \in D) = \frac{m}{n}, \quad \forall i \in \{1, \dots, n\}.
\end{equation}

\subsection{Algoritmo 1: Campionamento con Tabella Hash / Dizionario}
L'approccio più immediato prevede l'estrazione ripetuta di indici pseudocasuali compresi tra $1$ ed $n$, scartando i duplicati mediante una struttura dati per insiemi (\textit{Hash Set} o dizionario hash).

\begin{algorithm}[H]
\caption{Campionamento Casuale con Tabella Hash (\textit{HashSampling})}
\label{alg:hash-sampling}
\KwInput{Intervallo di indici $[1 \dots n]$, dimensione del campione $m \le n$}
\KwOutput{Insieme ordinato di $m$ indici distinti $D$}

Inizializza una tabella hash vuota $D \leftarrow \emptyset$\;
\While{$|D| < m$}{
  $p \leftarrow \text{rand}(1, n)$ \tcp*{Estrae intero uniforme in $[1, n]$}
  \If{$p \notin D$}{
    Inserisci $p$ in $D$\;
  }
}
Ordina gli elementi di $D$\;
\Return $D$\;
\end{algorithm}

\subsection{Analisi della Probabilità di Collisione e Tentativi Attesi}
A ogni iterazione del ciclo \texttt{while}, la probabilità che il numero casuale estratto $p$ appartenga già a $D$ (provocando una collisione e quindi un'estrazione nulla a vuoto) dipende dalla cardinalità corrente dell'insieme:
\begin{equation}
  \mathbb{P}(\text{collisione al passo con } |D| = j) = \frac{j}{n}.
\end{equation}

La probabilità di collisione raggiunge il proprio picco massimo nell'ultimo inserimento, quando l'insieme ha raggiunto cardinalità $m - 1$:
\begin{equation}
  \mathbb{P}(\text{collisione}) \le \frac{m - 1}{n} < \frac{m}{n}.
\end{equation}

\begin{property}[Efficienza per Campioni Ridotti ($m \le n/2$)]
Se la taglia del campione target soddisfa la condizione comune $m \le \frac{n}{2}$ (come nel caso dell'over-sampling, dove $m = \mathcal{O}(M/B \ln(M/B)) \ll n$), la probabilità di collisione a ogni iterazione è strettamente dominata da:
\begin{equation}
  \mathbb{P}(\text{collisione}) < \frac{1}{2}.
\end{equation}
Di conseguenza, più del $50\%$ dei numeri casuali generati corrisponde a un inserimento effettivo.
\end{property}

Il numero totale atteso di generazioni casuali per raccogliere tutti gli $m$ elementi distinti può essere espresso formalmente come somma di variabili geometriche indipendenti:
\begin{equation}
  \mathbb{E}[\text{tentativi}] = \sum_{j=0}^{m-1} \frac{1}{1 - \frac{j}{n}} = \sum_{j=0}^{m-1} \frac{n}{n - j} = n \sum_{i=n-m+1}^n \frac{1}{i}.
\end{equation}
Esprimendo la sommatoria mediante la differenza tra due numeri armonici:
\begin{equation}
  \mathbb{E}[\text{tentativi}] = n \left( H_n - H_{n-m} \right) \approx n \ln\left( \frac{n}{n - m} \right).
  \label{eq:sampling-expected-attempts}
\end{equation}

Effettuando lo sviluppo di Taylor di $\ln(1 + x)$ per $m \ll n$:
\begin{equation}
  n \ln\left(\frac{1}{1 - m/n}\right) = n \left( \frac{m}{n} + \frac{m^2}{2n^2} + \mathcal{O}\left(\frac{m^3}{n^3}\right) \right) = m + \frac{m^2}{2n} + \mathcal{O}\left(\frac{m^3}{n^2}\right).
\end{equation}
Quando il campione $m$ è contenuto rispetto a $n$ ($m = o(\sqrt{n})$), il numero di tentativi a vuoto è asintoticamente trascurabile e il costo temporale medio coincide con $\mathcal{O}(m)$.

\subsection{Limiti Strutturali e Modelli di Accesso}
Nonostante la semplicità implementativa, l'\autoref{alg:hash-sampling} presenta limiti architetturali rilevanti:

\begin{enumerate}
  \item \textbf{Accesso Indicizzato a Stringhe e Record a Lunghezza Variabile:} L'estrazione diretta dell'elemento $S[p]$ presuppone la capacità di indirizzare la cella $p$ in tempo $\mathcal{O}(1)$. Nei database e nei file di testo formattati contenenti stringhe o chiavi a lunghezza variabile, la posizione scalare di un record nel flusso di byte non è nota a priori senza aver scandito linearmente i record precedenti. L'estrazione casuale richiede pertanto un array ausiliario di puntatori o offset di memoria ($P[1 \dots n]$), introducendo un overhead di spazio $\Theta(n)$.
  
  \item \textbf{Assenza della Proprietà Streaming (\textit{One-Pass}):} L'algoritmo non è utilizzabile su flussi continui di dati (\textit{Data Streams}), poiché:
  \begin{itemize}
    \item Presuppone la conoscenza preventiva della cardinalità totale $n$.
    \item Richiede accesso casuale arbitrario bidirezionale all'intera sequenza di input.
  \end{itemize}
\end{enumerate}

\subsection{Campionamento su Flussi Continui: Il Reservoir Sampling}
Quando i dati transitano in una sequenza non memorizzabile per intero e la cardinalità $n$ non è nota a priori (o è potenzialmente infinita), si adotta l'algoritmo di \textbf{Reservoir Sampling} (Algoritmo R di Alan G. Waterman, formalizzato da Vitter~\cite{vitter1985}).

\begin{algorithm}[H]
\caption{Reservoir Sampling (\textit{Algoritmo R di Vitter})}
\label{alg:reservoir-sampling}
\KwInput{Flusso continuo di dati $S = \langle x_1, x_2, \dots \rangle$, dimensione serbatoio $m$}
\KwOutput{Serbatoio $R[1 \dots m]$ contenente $m$ campioni uniformi}

Inizializza il serbatoio con i primi $m$ elementi del flusso: $R[1 \dots m] \leftarrow [x_1, \dots, x_m]$\;
$t \leftarrow m$\;
\While{il flusso produce un nuovo elemento $x_{t+1}$}{
  $t \leftarrow t + 1$\;
  $j \leftarrow \text{rand}(1, t)$ \tcp*{Estrae intero uniforme in $[1, t]$}
  \If{$j \le m$}{
    $R[j] \leftarrow x_t$ \tcp*{Sostituisce il vecchio elemento in posizione $j$}
  }
}
\Return $R$\;
\end{algorithm}

\begin{theorem}[Correttezza del Reservoir Sampling]
Al termine della lettura di un qualsiasi prefisso di $n \ge m$ elementi del flusso, ogni elemento esaminato $x_i$ ($1 \le i \le n$) appartiene al serbatoio $R$ con probabilità rigorosamente uniforme e identica pari a:
\begin{equation}
  \mathbb{P}(x_i \in R) = \frac{m}{n}.
\end{equation}
\end{theorem}

\begin{proof}
Procediamo per induzione matematica sulla lunghezza del flusso $n \ge m$.
\begin{itemize}
  \item \textbf{Caso base ($n = m$):} I primi $m$ elementi vengono inseriti incondizionatamente in $R$. La probabilità per ciascuno di essi è $m/m = 1$, verificando la base induttiva.
  
  \item \textbf{Passo induttivo ($n > m$):} Assumiamo per ipotesi induttiva che, dopo aver esaminato $n - 1$ elementi, ogni elemento $x_i$ ($i \le n - 1$) risieda in $R$ con probabilità:
  \begin{equation}
    \mathbb{P}(x_i \in R \text{ al passo } n - 1) = \frac{m}{n - 1}.
  \end{equation}
  Consideriamo l'arrivo dell'$n$-esimo elemento $x_n$:
  \begin{enumerate}
    \item L'elemento $x_n$ viene inserito nel serbatoio se e solo se l'estrazione $j = \text{rand}(1, n)$ produce un valore $j \le m$. La probabilità di inserimento di $x_n$ è dunque:
    \begin{equation}
      \mathbb{P}(x_n \in R) = \frac{m}{n}.
    \end{equation}
    \item Un elemento generico $x_i$ preesistente in $R$ rimane nel serbatoio se non viene sostituito da $x_n$. La sostituzione si verifica con probabilità $\frac{m}{n} \cdot \frac{1}{m} = \frac{1}{n}$. Di conseguenza, la probabilità che $x_i$ sopravviva all'arrivo di $x_n$ è:
    \begin{equation}
      \mathbb{P}(x_i \text{ non rimpiazzato} \mid x_i \in R) = 1 - \frac{1}{n} = \frac{n - 1}{n}.
    \end{equation}
    Moltiplicando per la probabilità induttiva al passo precedente:
    \begin{equation}
      \mathbb{P}(x_i \in R \text{ al passo } n) = \frac{m}{n - 1} \cdot \frac{n - 1}{n} = \frac{m}{n}.
    \end{equation}
  \end{enumerate}
\end{itemize}
La proprietà è dimostrata per ogni $n \ge m$.
\end{proof}

\section{Sintesi Comparativa dei Risultati}
\label{sec:distribution-sort-synthesis}

A conclusione del capitolo, la \autoref{tab:sorting-paradigms-comparison} riassume il confronto metodologico e architetturale tra i due paradigmi d'ordinamento per memorie esterne, mentre la \autoref{tab:sampling-strategies-comparison} mette a confronto le strategie di campionamento casuale.

\begin{table}[H]
  \centering
  \caption{Confronto architetturale tra i macro-paradigmi di ordinamento su memoria secondaria.}
  \label{tab:sorting-paradigms-comparison}
  \begin{tabularx}{\textwidth}{@{}>{\bfseries\arraybackslash}p{0.22\textwidth}p{0.36\textwidth}p{0.36\textwidth}@{}}
    \toprule
    Caratteristica & Multiway Merge Sort (Cap. 3) & External Distribution Sort \\
    \midrule
    Filosofia & \textit{Bottom-Up} (fusione progressiva) & \textit{Top-Down} (partizionamento ricorsivo) \\
    \addlinespace
    Complessità I/O & $\mathcal{O}\left(\frac{N}{B} \log_{M/B} \frac{N}{B}\right)$ & $\mathcal{O}\left(\frac{N}{B} \log_{M/B} \frac{N}{B}\right)$ (in media/alta prob.) \\
    \addlinespace
    Buffer RAM & $k$ buffer di lettura, 1 buffer di scrittura & 1 buffer di lettura, $k$ buffer di scrittura \\
    \addlinespace
    Fase Terminale & Fusione complessa con Min-Heap & Concatenazione diretta dei secchi (0 I/O) \\
    \addlinespace
    Criticità & Elevato traffico I/O nella fase iniziale & Sensibilità allo sbilanciamento dei secchi \\
    \addlinespace
    Soluzione & Generazione run di taglia $2M$ (Loser Tree) & Over-Sampling ($a = 12 \ln k$) con Chernoff \\
    \bottomrule
  \end{tabularx}
\end{table}

\begin{table}[H]
  \centering
  \caption{Confronto analitico tra le strategie di campionamento casuale uniforme.}
  \label{tab:sampling-strategies-comparison}
  \begin{tabularx}{\textwidth}{@{}>{\bfseries\arraybackslash}p{0.22\textwidth}p{0.22\textwidth}p{0.24\textwidth}Y@{}}
    \toprule
    Algoritmo & Modello Dati & Costo Temporale & Proprietà e Limitazioni \\
    \midrule
    Tabella Hash (\textit{Rejection}) & Accesso Casuale $\mathcal{O}(1)$ su array & $\mathcal{O}(m)$ atteso (se $m \le n/2$) & Richiede conoscenza di $n$; inefficace su record a lunghezza variabile privi di indice. \\
    \addlinespace
    Reservoir Sampling (Vitter) & Flusso Streaming (\textit{One-Pass}) & $\mathcal{O}(n)$ & Non richiede di conoscere $n$; memoria limitata a $\mathcal{O}(m)$; legge i dati in sequenza stretta. \\
    \addlinespace
    Fisher-Yates Parziale & In-Place su array modificabile & $\mathcal{O}(m)$ & Modifica l'array originale scambiando le prime $m$ posizioni; richiede memoria scrivibile. \\
    \bottomrule
  \end{tabularx}
\end{table}
